\documentclass[aspectratio=169,11pt]{beamer}

\usetheme{Madrid}
\usecolortheme{default}
\usefonttheme{professionalfonts}
\setbeamertemplate{navigation symbols}{}
\setbeamertemplate{footline}[frame number]

\usepackage{amsmath, amssymb, booktabs, array, tabularx}

\title{Firm, Market, and Equilibrium Responses to Behavioral Frictions}
\subtitle{Behavioral Labor -- Week 8}
\author{Teaching deck draft}
\date{}

\newcolumntype{Y}{>{\raggedright\arraybackslash}X}

\begin{document}

\begin{frame}
  \titlepage
\end{frame}

\begin{frame}{Central question}
\begin{itemize}
    \item Once worker-level behavioral wedges exist, how do firms and markets respond?
    \item Week 8 follows the course sequence:
    \begin{itemize}
        \item Weeks 2--4: worker wedges in preferences, beliefs, attention, and complexity.
        \item Weeks 5--6: workplace and allocation environments.
        \item Week 7: measurement and identification.
    \end{itemize}
    \item Today: what survives after firm strategy, search, sorting, and competition?
\end{itemize}
\end{frame}

\begin{frame}{Worker-level wedges are not enough}
\[
Y_{im}=g\!\left(\theta_i,\;R_f(\theta_i,\mathcal{M}_m),\;\mathcal{M}_m\right)
\]

\begin{itemize}
    \item $\theta_i$: worker-level wedge such as inertia, misperception, or low salience.
    \item $R_f(\cdot)$: firm response through menus, defaults, messages, fees, wages, or recruiting.
    \item $\mathcal{M}_m$: market environment: competition, transparency, switching costs, search frictions, heterogeneity, market power.
    \item Observed choices are already equilibrium-shaped.
\end{itemize}
\end{frame}

\begin{frame}{Taxonomy of firm responses}
\scriptsize
\begin{tabularx}{\linewidth}{p{0.18\linewidth} Y Y}
\toprule
Response & Firm action & Empirical objects \\
\midrule
Exploit & design fees, defaults, or menus around inertia and confusion & shrouded attributes, switching, fee sensitivity, passive choice \\
Accommodate & simplify choices or make valuable actions easier & active choice, clearer menus, reduced mistakes, better matching \\
Insure & use defaults, restricted menus, or benefits to protect workers & contribution floors, benefit design, risk exposure, retention \\
Screen / sort & use messages, contracts, or menus to separate worker types & applicant pools, wage bids, plan choice by type, match quality \\
\bottomrule
\end{tabularx}

\vspace{0.6em}
The same practice can insure some workers and exploit or segment others.
\end{frame}

\begin{frame}{Employer-sponsored benefits as market response}
\begin{itemize}
    \item Benefits are compensation, but they are delivered through menus, defaults, deadlines, fees, and employer interfaces.
    \item Worker margin: enrollment, contribution, plan choice, switching.
    \item Firm or intermediary margin: menu design, default choice, disclosure, sales effort.
    \item Market margin: fee competition, demand elasticity, sorting, welfare.
    \item Core anchors: Bernheim--Fradkin--Popov; Handel--Kolstad; Duarte--Hastings; Hastings--Hortacsu--Syverson.
\end{itemize}
\end{frame}

\begin{frame}{401(k) defaults and welfare}
\small
\begin{tabularx}{\linewidth}{p{0.27\linewidth} Y}
\toprule
Question & Bernheim--Fradkin--Popov logic \\
\midrule
Identifying variation & workers face different retirement-plan defaults and menu environments \\
Observed margin & participation, contribution rates, asset allocation, default persistence \\
Worker-level effect & inertia, procrastination, passive choice \\
Firm response margin & employer default and menu design \\
Welfare issue & default-induced saving is not automatically welfare-improving \\
\bottomrule
\end{tabularx}

\vspace{0.6em}
Large default effects are evidence of a behavioral wedge; they are not a complete welfare statement.
\end{frame}

\begin{frame}{Plan choice, inertia, and transparency}
\scriptsize
\begin{tabularx}{\linewidth}{p{0.28\linewidth} Y Y}
\toprule
Paper domain & Observed margin & Equilibrium lesson \\
\midrule
Health insurance for ``humans'' & plan choice, expected costs, risk exposure & information frictions create welfare losses inside employer-mediated menus \\
Mexico social security disclosure & switching, fee sensitivity, account choice & transparency can raise demand elasticity, but firms may adjust marketing and sales effort \\
Sales force and competition & demand, fees, switching, plan shares & competition works through sales-force strategy, disclosure, and consumer sophistication \\
\bottomrule
\end{tabularx}

\vspace{0.6em}
The question is not only whether workers make mistakes. It is whether markets discipline or monetize those mistakes.
\end{frame}

\begin{frame}{Distorted outside-option beliefs}
\[
\tilde w_i^{out}=w_i^{out}-\mu_i
\]

\begin{itemize}
    \item Jager--Roth--Roussille--Schoefer: workers may misperceive outside options.
    \item Observed margins: beliefs, search, bargaining, mobility, wage response.
    \item If $\mu_i>0$, workers perceive outside options as worse than they are.
    \item Firms may retain workers at lower wages when perceived mobility is low.
    \item The belief object becomes a segmentation and market-power object.
\end{itemize}
\end{frame}

\begin{frame}{Search equilibrium and segmentation}
\begin{itemize}
    \item Week 3 question: do beliefs affect worker search?
    \item Week 8 question: do belief distortions change market allocation?
    \item Menzio's stubborn-beliefs logic asks when market experience fails to correct beliefs.
    \item Search frictions, weak wage transparency, and low mobility can preserve distorted beliefs.
    \item Frontier: behavioral monopsony, segmentation, and wage markdowns under distorted outside options.
\end{itemize}
\end{frame}

\begin{frame}{Firm messages, sorting, and search direction}
\begin{itemize}
    \item Horton--Kuchler--Stanton: employer messages can affect applications, wage bids, and sorting.
    \item Communication may reduce uncertainty and improve match quality.
    \item Communication may also screen worker types or steer search strategically.
    \item The empirical distinction:
    \begin{itemize}
        \item worker response to messages;
        \item firm choice of messages;
        \item market effect on matching and wage dispersion.
    \end{itemize}
\end{itemize}
\end{frame}

\begin{frame}{When do wedges survive?}
\scriptsize
\begin{tabularx}{\linewidth}{p{0.23\linewidth} Y Y}
\toprule
Condition & Attenuation & Survival or amplification \\
\midrule
Competition & clear total compensation and better menus & competition on salient attributes while shrouding hidden costs \\
Transparency & workers compare wages, benefits, fees, outside options & attributes are opaque, delayed, or broker mediated \\
Switching costs & workers can change plans or jobs cheaply & default persistence and slow job or plan switching \\
Search frictions & offers and outside options are learned quickly & low search and segmented information networks \\
Heterogeneity & sophisticated workers discipline markets & firms target less sophisticated workers \\
Labor market power & workers can credibly leave & firms retain workers with pessimistic beliefs or low mobility \\
\bottomrule
\end{tabularx}
\end{frame}

\begin{frame}{Identification map}
\scriptsize
\begin{tabularx}{\linewidth}{p{0.25\linewidth} Y Y}
\toprule
Block & Identifying variation & First object identified \\
\midrule
Retirement defaults & default exposure and menu variation & worker default response under a given menu \\
Health-plan choice & information and plan-choice environment & plan-choice mistakes and welfare losses \\
Retirement account markets & disclosure, sales force, competition & demand elasticity, switching, fee sensitivity \\
Outside-option beliefs & elicited beliefs and objective outside-option measures & belief distortion and worker search response \\
Labor-market messages & message or platform variation & application, bid, and sorting response \\
\bottomrule
\end{tabularx}

\vspace{0.5em}
Equilibrium interpretation needs extra evidence on firm response and market outcomes.
\end{frame}

\begin{frame}{Partial versus equilibrium welfare}
\begin{itemize}
    \item Attenuation is not always correction.
    \item Amplification is not always manipulation.
    \item A partial effect can disappear or reverse after firms redesign menus, fees, wages, or messages.
    \item Revealed preference is weaker when choice architecture is endogenous.
    \item Welfare requires a benchmark: full information, low hassle, sophisticated choice, structural utility, or planner target.
\end{itemize}
\end{frame}

\begin{frame}{Research lab}
\begin{itemize}
    \item Reproduce: retirement-default exercise anchored on Bernheim--Fradkin--Popov.
    \item Diagnose: menu design, default quality, passive choice, and welfare distance.
    \item Transfer: outside-option beliefs and segmentation anchored on Jager--Roth--Roussille--Schoefer.
    \item Optional extension: disclosure, sales effort, demand elasticity, and switching costs from Duarte--Hastings or Hastings--Hortacsu--Syverson.
\end{itemize}
\end{frame}

\begin{frame}{Bridge to Week 9 policy}
\begin{itemize}
    \item Week 9 asks how policy should be designed once firms and markets respond.
    \item Defaults, disclosure, simplification, and nudges are not implemented in a vacuum.
    \item Policy can change worker behavior and firm strategy at the same time.
    \item The policy target must name the worker wedge, firm response, market condition, and welfare benchmark.
\end{itemize}
\end{frame}

\begin{frame}{Takeaways}
\begin{itemize}
    \item Worker-level behavioral wedges are not observed market outcomes.
    \item Firms can exploit, accommodate, insure, screen, sort, and segment.
    \item Competition helps only under the right transparency, switching, and search conditions.
    \item The empirical frontier links worker behavior to firm response and market outcomes.
    \item Welfare and policy require equilibrium interpretation, not only strong reduced forms.
\end{itemize}
\end{frame}

\end{document}
